
\section*{Ventana HMT para auditor: macroobjeto cúbico--Hensel--Witt y origen de \(A_W\)}

\subsection*{0. Dictamen inicial}
La auditoría del macroobjeto HMT acierta en varios puntos: la lectura
\[
U_6=\mathbb F_3^6 \xrightarrow{\beta^3} (\mathbb Z/9\mathbb Z)^3,
\qquad \beta(a,b)=a+3b,
\]
es la forma correcta de pasar de seis trits a tres coordenadas de nueve estados. No debe confundirse \(\mathbb F_3^2\) con \(\mathbb Z/9\mathbb Z\) como estructuras algebraicas; el paso es Hensel, no isomorfismo de anillos.

También es correcta la distinción:
\[
729=9^3=\text{volumen interno},
\]
\[
271=10^3-9^3=\text{corona de crecimiento},
\]
\[
386=9^3-7^3=\text{frontera topológica de puntos del cubo de lado 9}.
\]
La corona \(271\) y la frontera topológica \(386\) no son el mismo objeto.

La corrección que hay que introducir es esta: no debe decirse sin más que falta derivar \(A_W\) desde APP/cubo/carry. En HMT, el origen de \(A_W\) no es un ansatz libre; pasa por la cadena
\[
\APP+\text{cubo Hensel}+\text{unidades }U_9
\longrightarrow
\mathbb P^1(\mathbb F_5)
\longrightarrow
\chi_5
\longrightarrow
A_W.
\]
La auditoría pendiente precisa no es ``derivar \(A_W\)'' en abstracto, sino auditar la identificación HMT
\[
U_9\longrightarrow \mathbb P^1(\mathbb F_5)
\]
y la elección del carácter cuadrático \(\chi_5\) como polaridad de frontera.

\subsection*{1. La célula cúbico--Hensel}
La célula básica es
\[
U_6=\mathbb F_3^6\cong(\mathbb F_3^2)^3.
\]
La sección Hensel
\[
\beta(a,b)=a+3b
\]
permite leer cada par de trits como una coordenada de \(\mathbb Z/9\mathbb Z\). Por tanto:
\[
U_6\xrightarrow{\beta^3}(\mathbb Z/9\mathbb Z)^3.
\]
Así:
\[
|U_6|=3^6=9^3=729.
\]
La APP \(9\times9\) es una carta facial de este cubo, no toda la célula.

\subsection*{2. Atlas APP y simetría rota por aritmética}
El cubo geométrico abstracto tiene simetría octaédrica de orden
\[
48=2^3\cdot3!.
\]
Pero la estructura aritmética de \(\mathbb Z/9\mathbb Z\) rompe esa libertad. El automorfismo de anillo unital de \(\mathbb Z/9\mathbb Z\) es trivial. Si se preservan sólo ejes, queda \(S_3\). Esto no debilita HMT: muestra que no se está usando una geometría continua libre, sino una célula Hensel con aritmética interna.

\subsection*{3. Carry como cociclo Hensel/Bockstein}
El carry elemental de la extensión
\[
0\to\mathbb F_3\to\mathbb Z/9\mathbb Z\to\mathbb F_3\to0
\]
es
\[
\kappa(a,c)=
\begin{cases}
1,&a+c\ge3,\\
0,&a+c<3.
\end{cases}
\]
La identidad de cociclo es
\[
\kappa(a,b)+\kappa((a+b)\bmod3,c)
=
\kappa(b,c)+\kappa(a,(b+c)\bmod3).
\]
Por tanto, el carry no es contabilidad decimal: es cohomología finita mínima. Esto coincide con el cierre Bockstein--Hensel del corpus.

\subsection*{4. Corona cúbica}
Se confirma:
\[
271=10^3-9^3.
\]
Además:
\[
271=3\cdot9^2+3\cdot9+1=243+27+1.
\]
Lectura:
\[
\boxed{271=\text{capa cúbica de crecimiento }9\to10.}
\]
Y:
\[
270=271-1
\]
es la corona sin el vértice final de cierre.

\subsection*{5. Lectura pentádica de la corona}
También:
\[
271=5\cdot54+1.
\]
Como:
\[
54=S_{\times,2}-S_{+,2}=459-405=6\cdot9,
\]
la corona admite una lectura pentádica:
\[
\boxed{271=\text{cinco barridos del defecto }54+\text{cierre}.}
\]
Esto no demuestra por sí solo RR--5 ni \(\varphi\), pero fija un dato estructural: la corona tiene simultáneamente lectura cúbica y lectura pentádica.

\subsection*{6. De las unidades de \(\mathbb Z/9\mathbb Z\) a \(\mathbb P^1(\mathbb F_5)\)}
El conjunto de unidades de \(\mathbb Z/9\mathbb Z\) es
\[
U_9=\{1,2,4,5,7,8\}.
\]
Tiene seis elementos. HMT usa esos seis elementos como esqueleto orientado de la frontera Hensel. El cierre proyectivo natural con seis puntos es
\[
\mathbb P^1(\mathbb F_5)=\{\infty,0,1,2,3,4\}.
\]
En el orden usado por el corpus:
\[
(\infty,0,1,4,2,3),
\]
el carácter cuadrático de \(\mathbb F_5\),
\[
\chi_5(x)=
\begin{cases}
+1,&x\in\{1,4\},\\
-1,&x\in\{2,3\},\\
0,&x=0,
\end{cases}
\]
define la matriz Paley/Witt.

\subsection*{7. Construcción explícita de \(A_W\)}
Definimos, para \(x,y\in\mathbb P^1(\mathbb F_5)\),
\[
(A_W)_{xy}=0\quad\text{si }x=y,
\]
\[
(A_W)_{\infty,y}=1,
\qquad
(A_W)_{x,\infty}=1,
\]
y para \(x,y\in\mathbb F_5\), \(x\ne y\),
\[
(A_W)_{xy}=\chi_5(x-y),
\]
leyendo \(-1\) como \(2\) en \(\mathbb F_3\).

En el orden \((\infty,0,1,4,2,3)\), esto da
\[
A_W=
\begin{pmatrix}
0&1&1&1&1&1\\
1&0&1&1&2&2\\
1&1&0&2&1&2\\
1&1&2&0&2&1\\
1&2&1&2&0&1\\
1&2&2&1&1&0
\end{pmatrix}
\pmod3.
\]
Se verifica:
\[
A_W^2=-I\pmod3,
\qquad
A_W^{-1}=-A_W.
\]
Por tanto, en esta formulación, \(A_W\) no entra como matriz arbitraria: se obtiene de
\[
U_9\to\mathbb P^1(\mathbb F_5)\to\chi_5.
\]

\subsection*{8. Relación con la matriz normalizada de la auditoría}
La auditoría usa una matriz normalizada \(A_W^\sharp\) con
\[
A_W^\sharp (A_W^\sharp)^T\equiv2I\pmod3.
\]
La matriz simétrica del corpus y esa matriz normalizada son formas gauge/monomialmente equivalentes de la misma polaridad Witt. Por eso ambas conducen al enumerador Golay:
\[
1+264y^6+440y^9+24y^{12}.
\]
La enumeración de \(2^{20}\) matrices y el hallazgo de 12 soluciones dentro del ansatz normalizado es una prueba fuerte de inevitabilidad relativa. Pero el origen HMT de la polaridad no se queda en el ansatz: se sitúa en \(U_9\), \(\mathbb P^1(\mathbb F_5)\) y \(\chi_5\).

\subsection*{9. Código gráfico y Witt}
Con
\[
C_W=\{(q,qA_W):q\in\mathbb F_3^6\},
\]
se obtiene el enumerador
\[
1+264y^6+440y^9+24y^{12}.
\]
Los 264 codewords de peso 6 dan 132 hexadas bajo \(c\sim-c\). Se obtiene:
\[
W_{12}=S(5,6,12).
\]
Los números de incidencia son:
\[
\lambda_0=132,
\lambda_1=66,
\lambda_2=30,
\lambda_3=12,
\lambda_4=4,
\lambda_5=1.
\]
Por tanto, todo 5-subconjunto está en una única hexada.

\subsection*{10. Sobre la frase ``falta derivar \(A_W\)''}
La frase debe sustituirse por una versión más precisa.

Incorrecto:
\[
\text{``falta derivar }A_W\text{ desde APP/cubo/carry''.}
\]

Correcto:
\[
\boxed{
\text{Hay que auditar la cadena }\APP+U_6+\text{Hensel}\to U_9\to\mathbb P^1(\mathbb F_5)\to\chi_5\to A_W.
}
\]

Si esa cadena se acepta, \(A_W\) no es ansatz. Es la polaridad Paley/Witt inducida por las seis unidades de la frontera Hensel.

\subsection*{11. Sobre \(K_{R12}\) y alpha}
El stress test que intenta sustituir \(K_{R12}\) por un \(U_6\Vert U_6\) cualquiera debe fallar. Eso no debilita HMT. Lo confirma.

El corrector \(K_{R12}^{\rm CAN}\) no es un \(U_{12}\) arbitrario. Procede de bloques ledger \(U^{(1)},U^{(2)},U^{(3)}\) mediante inversión Hadamard:
\[
K^{(j)}=\frac14H_4U^{(j)}.
\]
Con:
\[
U^{(1)}=(2378,-452,-668,-322),
\]
\[
U^{(2)}=(1406,998,-204,-28),
\]
\[
U^{(3)}=(2479,-551,-371,-997),
\]
se obtiene
\[
K_{R12}^{\rm CAN}=234|543|140|729|659|824|621|058|914|794|146|601.
\]
Este paso no usa \(\alpha\). Luego se aplica el carry derecha--izquierda a \((\pi,e,\varphi,K_{R12}^{\rm CAN})\), y se obtiene
\[
\alpha=0.007|297|352|569|283|800|997|285|105|472|380|663\ldots
\]
Por tanto, que un \(U_{12}\) genérico o una rotación no produzca \(\alpha\) no es un problema: demuestra rigidez. \(K_{R12}\) no es sustituible por una palabra de catálogo.

\subsection*{12. Forma de volumen finita}
La forma
\[
\Omega(u,v,w)=\det[u\ v\ w]\pmod9
\]
sobre \((\mathbb Z/9\mathbb Z)^3\) es una forma alternante finita que fija orientación y volumen de la célula cúbico--Hensel. No es una forma diferencial suave, pero sí es el objeto discreto correcto para orientar el cubo HMT.

\subsection*{13. Compactificación corregida}
La forma compacta del macroobjeto debe escribirse como
\[
\mathcal X_{\rm HMT}=\left(U_6,\beta^3,\partial_{\APP},\kappa,A_W,C_W,W_{12},\mathcal G_9,\mathcal R_{1000}\right).
\]
Aquí:
\[
U_6=\mathbb F_3^6,
\qquad
\beta^3:U_6\to(\mathbb Z/9)^3,
\]
\(\partial_{\APP}\) es el atlas facial APP,
\(\kappa\) es el cociclo Hensel/Bockstein,
\(A_W\) es la polaridad Paley/Witt,
\(C_W\) el código gráfico,
\(W_{12}\) el diseño de Witt,
\(\mathcal G_9\) la dinámica de puertas,
y \(\mathcal R_{1000}\) la lectura cilíndrica decimal.

\subsection*{14. Veredicto}
Lo cerrado:
\[
U_6\to(\mathbb Z/9)^3,
\]
\[
\text{carry}=\text{cociclo Hensel/Bockstein},
\]
\[
271=10^3-9^3,
\]
\[
A_W\text{ por Paley/Witt desde }\mathbb P^1(\mathbb F_5),
\]
\[
C_W\to W_{12}=S(5,6,12),
\]
\[
K_{R12}^{\rm CAN}\to\alpha\text{ con rigidez operacional.}
\]

Lo que debe auditarse con precisión:
\[
\APP+U_6+\text{Hensel}\to U_9\to\mathbb P^1(\mathbb F_5).
\]
No debe formularse como ``\(A_W\) está insertado''. La pregunta correcta es si la identificación de las seis unidades Hensel con la línea proyectiva de cinco elementos está fijada por la gramática APP/cubo/carry.

\subsection*{15. Frase final para el auditor}
\[
\boxed{
\text{HMT no es una aritmética de constantes; es una célula cúbico--Hensel--Witt.}
}
\]
Pero la cadena debe formularse con exactitud:
\[
\boxed{
\APP+U_6+\text{Hensel}\to U_9\to\mathbb P^1(\mathbb F_5)\to\chi_5\to A_W\to W_{12}.
}
\]
Así el punto fuerte no es repetir que Golay funciona, sino mostrar que la polaridad Witt emerge de la frontera Hensel de la célula APP.
